% Generated draft: requires mathematical and source review.
\documentclass[11pt]{article}
\usepackage[margin=25mm]{geometry}
\usepackage{fontspec}
\setmainfont{Latin Modern Roman}
\usepackage{amsmath,amssymb,mathtools,unicode-math}
\setmathfont{Latin Modern Math}
\usepackage{enumitem,needspace,xurl,hyperref}
\setlength{\parindent}{0pt}
\setlength{\parskip}{5pt}
\setlength{\emergencystretch}{3em}
\newcommand{\N}{\mathbb{N}}
\newcommand{\Z}{\mathbb{Z}}
\newcommand{\Q}{\mathbb{Q}}
\newcommand{\R}{\mathbb{R}}
\newcommand{\C}{\mathbb{C}}
\newcommand{\F}{\mathbb{F}}
\newcommand{\A}{\mathbb{A}}
\newcommand{\PP}{\mathbb{P}}
\newcommand{\E}{\mathbb{E}}
\newcommand{\eps}{\varepsilon}
\newcommand{\dd}{\,\mathrm{d}}
\begin{document}
\section*{UnsolvedMath ID: 2533}
\textbf{Source problem code:} catalogue entry
\subsection{Definitions and mathematical statement}
Let $G$ be a finite group and $\langle x\rangle=\{x^j:j\in\Z\}$ its cyclic subgroup. Its undirected power graph $P(G)$ is the finite simple graph with all elements of $G$ as vertices, including the identity, and
\[
\{x,y\}\in E(P(G))\quad\Longleftrightarrow\quad
x\ne y\text{ and }(x\in\langle y\rangle\text{ or }y\in\langle x\rangle).
\]
A cograph is a graph with no induced path on four vertices. A chordal graph has no induced cycle of length at least four; equivalently every cycle of length at least four has a chord. The target is
\[
\forall G\text{ finite group}:\qquad
P(G)\text{ a cograph}\ \Longrightarrow\ P(G)\text{ chordal}.
\]
The graph is the full power graph, not the identity-deleted proper power graph, and no simplicity or element-order restriction is imposed on $G$.

The supplied source is Kourovka Notebook, New Problems, 21st issue (2026), Problem 21.24, cross-referenced to 21.23. The background names D. Bubboloni, F. Fumagalli and C. E. Praeger and reports earlier prime-power-element-order and nonabelian-simple-group cases. It then gives independent affirmative manuscripts of Rundström dated 30 January 2026 and Monticone dated 29 March 2026.

The literature triage, checked 2026-08-17, reports the implication proved for every finite group and accompanied by a Lean formalization in a 2026 Kourovka project. It labels the question solved with no remaining mathematical component, while journal publication of recent manuscripts remains pending. The primary registry status is also ``solved''. These are supplied proof-status reports, not an independent proof or formal-verification run.
\subsection{Short English statement}
If a finite group's full power graph has no induced four-vertex path, must it also have no induced cycle of length at least four? The supplied 2026 record reports an affirmative theorem for every finite group and an accompanying Lean formalization.
\subsection{Sources}
\begin{itemize}
\item UnsolvedMath local registry, permanent ID 2533, source code catalogue entry, ``Kourovka Notebook Problem 21.24'', New Problems, 21st issue (2026). The background names D. Bubboloni, F. Fumagalli and C. E. Praeger and says ``Cf. 21.23'' without that statement. Dataset attribution: Ulam AI and the UnsolvedMath Contributors; curation and metadata under CC BY 4.0.
\item D. Bubboloni, F. Fumagalli and C. E. Praeger, preprint (2025), title not supplied, \url{https://arxiv.org/abs/2510.18073}. The background cites the prime-power-order special case.
\item J. Cameron, P. Manna and R. Mehatari, Journal of Algebra 591 (2022), 59--74; J. Brachter and E. Kaja, Journal of Algebraic Combinatorics 58 (2023), 1095--1124. Titles and direct URLs are not supplied; the background cites nonabelian-simple-group cases.
\item M. Rundström, solution preprint for Problem 21.24, recorded as 30 January 2026, \url{https://kourovkanotebookorg.wordpress.com/wp-content/uploads/2026/04/21.24-runds.pdf}. The supplied review identifies an independent affirmative manuscript.
\item P. Monticone, preprint recorded as 29 March 2026, \url{https://kourovkanotebookorg.wordpress.com/wp-content/uploads/2026/04/21_24-1.pdf}, title not supplied.
\item Wouter van Doorn, Elias Judin, Pietro Monticone and Daniel Morrison, \emph{On Some Problems from the Kourovka Notebook}, arXiv:2607.17477 (2026), \url{https://arxiv.org/abs/2607.17477}.
\item Pietro Monticone et al., Kourovka formalization repository, Problem 21.24 (2026), \url{https://github.com/pitmonticone/Kourovka}. The supplied assessment cites its indexed Lean proof.
\end{itemize}
\end{document}
